\documentclass[11pt]{article}

\usepackage[margin=1in]{geometry}
\usepackage{amsmath,amssymb}
\usepackage{booktabs}
\usepackage{listings}
\usepackage{xcolor}
\usepackage{hyperref}
\usepackage[numbers,sort&compress]{natbib}

\hypersetup{colorlinks=true,linkcolor=blue!50!black,citecolor=blue!50!black,
            urlcolor=blue!50!black}

% Draft scaffolding.  Everything wrapped in \pending is an argument the paper
% intends to make and has not yet earned; \todo marks a smaller loose end.
% Both render loudly on purpose.  Remove this block before submission and the
% document should still compile.
\newcommand{\pending}[1]{%
  \par\medskip\noindent\fcolorbox{red!60!black}{red!4}{%
    \parbox{\dimexpr\linewidth-2\fboxsep-2\fboxrule}{\small\textbf{[PENDING]}\ #1}}%
  \par\medskip}
\newcommand{\todo}[1]{\textcolor{red!70!black}{\small[\textbf{TODO:} #1]}}

\lstset{basicstyle=\ttfamily\small,columns=fullflexible,keepspaces=true,
        showstringspaces=false,frame=single,framesep=4pt,rulecolor=\color{black!25}}

\title{Parsimony: Revisiting Search-Based Floating-Point\\
       Precision Tuning for Reduced-Precision Hardware}
\author{Mark Adams\\
        \small Lawrence Berkeley National Laboratory\\
        \small \texttt{mfadams@lbl.gov}}
\date{Draft --- \today}

\begin{document}
\maketitle

\begin{abstract}
Automated floating-point precision tuning was solved, in method, over a decade
ago. Precimonious~\cite{precimonious} searched type configurations with delta
debugging; Blame Analysis~\cite{blame} made the search tractable by pruning it
with a single shadow execution. Both are unmaintained. We argue this is not
because the methods failed but because the hardware of the time did not reward
them: in the Blame Analysis evaluation, four of ten benchmark programs show
exactly $0.0\%$ speedup at every error threshold tested, and a fifth never
exceeds $0.4\%$. Converting \texttt{double} to \texttt{float} on a 2013 CPU
bought a little vectorization and cache residency, net of the casts it
introduced, and the tools frequently reported that the original program was
fastest.

The objective function has since changed. Reduced precision now addresses tensor
cores, \texttt{bfloat16} and 8-bit formats, and bandwidth-bound kernels in which
halving a matrix's storage approaches a factor of two in time. We present
\emph{parsimony}, a precision-tuning framework that separates the two concerns
prior tools fused --- how a precision configuration is \emph{evaluated}, and how
the configuration space is \emph{searched} --- and reimplements the delta-debugging
search against modern formats. We validate the emulation against hardware and the
search against a synthetic oracle, and reproduce a published Blame Analysis result
at a finer granularity than the original tool could express.
\pending{Quantitative results on NAS \texttt{cg} (\S\ref{sec:nascg}) and on
PETSc solvers (\S\ref{sec:petsc}). The abstract should close with a headline
number from one of them.}
\end{abstract}

\section{Introduction}
\label{sec:intro}

Choosing floating-point precision is a decision most scientific software makes
once, conservatively, and never revisits. Declaring everything \texttt{double} is
defensible, cheap to reason about, and --- on the hardware most of this software
was written for --- nearly free. That last clause no longer holds.

Automated tuning tools for this problem appeared a decade ago and worked.
Precimonious~\cite{precimonious} treats each floating-point variable as taking a
type from $\{\texttt{float}, \texttt{double}, \texttt{long double}\}$ and searches
that space with Zeller and Hildebrandt's delta-debugging
algorithm~\cite{ddmin}, seeking an assignment that both meets an accuracy
threshold and runs faster. Blame Analysis~\cite{blame} attacks the resulting
scalability problem: rather than recompiling and rerunning per candidate, it
executes the program once under shadow execution, computing for each instruction
the minimum operand precision that still yields the required result precision,
and uses the result to prune Precimonious' search space --- $9\times$ faster
analysis on average, up to $38\times$.

Neither tool is maintained. The natural reading is that the approach did not pan
out. We think the evidence supports a different reading, and it is visible in the
tools' own numbers.

\subsection{The observation}

Table~5 of the Blame Analysis paper reports program speedup after tuning, across
ten programs and four error thresholds ($10^{-4}$ through $10^{-10}$), for both
Precimonious alone and Blame\,+\,Precimonious. For \texttt{bessel},
\texttt{gaussian}, \texttt{roots}, and \texttt{rootnewt}, every entry is
$0.0\%$. For \texttt{polyroots} no entry exceeds $0.4\%$. Real speedups appear
only for \texttt{sum} ($39.9\%$), \texttt{fft} ($8.3\%$), \texttt{blas}
($5.1\%$), NAS \texttt{cg} ($7.7$--$7.9\%$), and \texttt{ep} ($0.6\%$). In the
configuration tables, a majority of cells are marked to indicate that the tools
selected the \emph{original} program as fastest.
\todo{Verify each cited cell against the published table before submission; the
figures above were read from the paper's typeset tables.}

This is not a failure of search. The searches found valid lower-precision
configurations --- Blame Analysis reports that on average $40\%$ of program
variables can be declared \texttt{float} (median $28\%$). It is a failure of
\emph{payoff}. On a 2013 general-purpose CPU, both \texttt{float} and
\texttt{double} arithmetic issue at similar rates; the gains available are
narrower vector lanes and better cache residency, and they are offset by the
conversion instructions a mixed configuration introduces. The Blame Analysis
authors say so directly, noting that lowering precision does not necessarily
produce a faster program and that a triage by the programmer is required.

\subsection{Why this is now a different problem}

Three changes since 2016 alter the arithmetic of that trade-off.

\paragraph{Formats below \texttt{float}.} The tuning lattice of the original
tools was $\{\texttt{float}, \texttt{double}, \texttt{long double}\}$, with
intermediate \emph{virtual} precisions used internally by Blame Analysis but not
expressible as storage. Hardware now provides \texttt{bfloat16}, \texttt{fp16},
\texttt{tf32}, and two 8-bit encodings. A variable that Precimonious could at
best move from 64 to 32 bits may now move to 8, and the search space is
correspondingly richer.

\paragraph{Throughput that actually differs by format.} Matrix units on current
accelerators deliver order-of-magnitude higher throughput at reduced precision.
The performance term in the dual objective, which was near zero for most
programs in the original evaluation, is no longer near zero.

\paragraph{Bandwidth-bound kernels.} Much of scientific computing --- sparse
matrix-vector products above all --- is limited by memory traffic rather than by
arithmetic. For these, precision is a \emph{storage} decision, and halving the
width of the dominant array approaches a factor of two in time regardless of what
the arithmetic units do. This case barely existed in the original benchmark set,
which was drawn from GSL routines and small NAS kernels.

\subsection{Contributions}

\begin{itemize}
  \item An argument, grounded in the prior tools' published evaluations, that
    search-based precision tuning was limited by its era's hardware payoff rather
    than by its method (\S\ref{sec:background}).
  \item A framework, \emph{parsimony}, that separates evaluation from search as
    independent plugin axes, so that a cheap evaluator can drive an expensive
    search during development and an expensive evaluator confirm the result
    (\S\ref{sec:design}).
  \item Clean-room reimplementations of delta-debugging precision search
    generalized to arbitrary $(\text{exponent}, \text{significand})$ formats,
    validated against a synthetic oracle independently of any floating-point
    emulation (\S\ref{sec:search}, \S\ref{sec:validation}).
  \item Reduced-precision emulation verified bit-identical to hardware
    \texttt{float32} and \texttt{float16}, including subnormals
    (\S\ref{sec:emulation}).
  \item An independent reproduction of a published Blame Analysis result, at a
    granularity the original lattice could not express
    (\S\ref{sec:repro}).
\end{itemize}

\pending{A contribution bullet reporting a quantitative result on a real solver.
Until \S\ref{sec:petsc} exists this paper is an architecture and a reproduction,
which is not yet an arXiv-worthy claim.}

\section{Background}
\label{sec:background}

\subsection{Precimonious}

Precimonious~\cite{precimonious} operates on LLVM bitcode. Given a program, a set
of test inputs, and an error threshold, it searches for a \emph{type
configuration}: an assignment of a floating-point type to each program variable.
The search is delta debugging, with worst-case complexity $O(n^2)$ in the number
of variables, and each candidate requires recompilation and re-execution. The
paper reports cases in which analysis takes hours for programs declaring fewer
than 50 variables with native runtimes under five seconds.

Two practical constraints deserve note because they persist in any successor.
First, the analysis is performed without optimization, because optimization
eliminates the very variables the search ranges over --- so the configuration is
tuned on a program that differs from the one that ships. Second, the search is
empirical: a configuration meeting $10^{-6}$ will satisfy $10^{-4}$, but the
configuration found directly for $10^{-4}$ may differ, and may perform better.

\subsection{Blame Analysis}

Blame Analysis~\cite{blame} replaces search with attribution. Each floating-point
instruction is executed under shadow execution at several operand precisions
simultaneously --- \texttt{float}, \texttt{double}, and \texttt{double} truncated
to 4, 6, 8 and 10 significant decimal digits. For a given precision requirement
on an instruction's result, the analysis records the minimal operand precisions
that satisfy it; this \emph{blame set} is merged across dynamic instances of the
instruction by taking the maximum requirement. A backward pass from the target
statement then yields the set of variables that may be single precision.

The cost structure matters for anyone considering the approach at scale. The
online variant is bounded by the static instruction count and imposes up to
$50\times$ runtime overhead. The offline variant finds better configurations but
is bounded by the \emph{dynamic} trace: it exhausted 256\,GB of memory on NAS
\texttt{ep} at class~S, the smallest input designed for serial execution.

One remark in that paper is the germ of the present work's approach to scale: the
authors observe that long-running scientific programs generally use iterative
solvers, that precision requirements tighten toward the end of execution, and
that analyzing only the last few iterations is therefore likely to suffice. In
other words, deciding \emph{what} to analyze is the lever, and it is a modeling
decision rather than an algorithmic one.

\section{Design}
\label{sec:design}

\subsection{Two orthogonal axes}

Prior tools bind their search strategy to their measurement technology.
Precimonious' delta debugging is inseparable from its LLVM~3.0 instrumentation;
PROMISE's search is inseparable from CADNA. The consequence is that neither
survives its dependency, and that a search improvement cannot be evaluated
against a different error model.

Parsimony defines two protocols:

\begin{lstlisting}[language=Python]
class Evaluator(Protocol):
    caps: Caps
    def discover(self) -> ConfigSpace: ...
    def evaluate(self, config: Config) -> Result: ...

class Search(Protocol):
    def search(self, space: ConfigSpace,
               evaluate: Callable[[Config], Result]) -> dict[str, Format]: ...
\end{lstlisting}

An \texttt{Evaluator} answers what a configuration costs in accuracy, and
optionally in time; a \texttt{Search} decides which configurations to ask about.
\texttt{Caps} declares what an evaluator can observe --- which formats, whether it
requires a rebuild, whether it sees GPU code, its slowdown --- so that a request no
installed evaluator can serve fails explicitly rather than being silently served
by something weaker.

The separation has a methodological payoff beyond modularity: it makes the search
testable in isolation. A synthetic oracle with a known minimal set exercises the
search with no floating-point emulation involved at all
(\S\ref{sec:validation}), so a regression in the search is diagnosed where it
occurs rather than surfacing as an inexplicably poor precision assignment.

\subsection{Formats and emulation}
\label{sec:emulation}

A format is an (exponent bits, significand bits) pair, which covers the IEEE
formats, the vendor formats, and \emph{virtual} widths --- a \texttt{float64}
truncated to $k$ significand bits with the full exponent range. Virtual widths
are how one asks how many bits a value needs without committing to a storage
format, and they generalize the truncated-decimal lattice Blame Analysis used
internally.

Rounding is emulated on \texttt{float64}: the significand is rounded to nearest
via \texttt{frexp}/\texttt{ldexp} scaling, values above the format's largest
finite magnitude become infinite, and values below its minimum normal exponent
are quantized to the appropriate subnormal step. This is exact for any format
whose range and precision fit within \texttt{float64}, which is every format
below it; wider formats require a native or instrumented evaluator.

\subsection{Search}
\label{sec:search}

The delta-debugging adaptation requires one observation to state correctly.
Delta debugging minimizes a subset subject to a predicate that holds on the full
set. Here the subset is the set of variables \emph{retained at higher precision},
and the predicate is that the accuracy goal is met. That predicate is monotone in
the direction the algorithm assumes --- returning a variable to higher precision
cannot hurt accuracy --- so the standard algorithm applies unmodified. Where
monotonicity fails, as when a coarser format is accidentally favorable to a
particular cancellation, the algorithm still terminates and the result is locally
rather than globally minimal, which is equally true of the original tools.

The search walks the format ladder from high to low. At each rung, the variables
still above that rung are candidates for demotion; delta debugging identifies the
minimal set that must remain above it; the remainder are demoted and the search
proceeds. This costs $O(n^2)$ evaluations \emph{per rung}, which is the dominant
scalability concern and the motivation for the prefilter discussed in
\S\ref{sec:limitations}.

We also provide uniform bisection over the ladder, which finds the lowest single
format that suffices in $O(\log |\text{ladder}|)$ evaluations. This is not a
competitor to the mixed search but a control for it: the gap between the uniform
floor and the mixed configuration is the only honest measure of what the
expensive search bought.

\section{Validation}
\label{sec:validation}

\subsection{Emulation against hardware}

Emulated rounding to \texttt{fp32} and \texttt{fp16} is bit-identical to numpy's
native \texttt{float32} and \texttt{float16} over random samples, including
inputs in the subnormal range, and \texttt{bfloat16} is verified to retain
\texttt{fp32}'s exponent range on values that overflow \texttt{fp16}.

\subsection{Search against a synthetic oracle}

The delta-debugging implementation is tested against a predicate that holds
exactly when a known subset is present. The search recovers that subset exactly,
its output is verified 1-minimal, and the evaluation count is bounded well below
exhaustive enumeration. No floating-point emulation participates in these tests.

\subsection{Reproducing a published result}
\label{sec:repro}

The worked example of the Blame Analysis paper (its Figure~2) computes
$a^4 - 4a^3 + 6a^2 - 4a + 1$ at $a = 1.84089642$ by repeated multiplication.
Intermediate terms reach roughly $25$ while the result is near $0.5$, so the
summation loses about $1.7$ decimal digits to cancellation. The paper concludes
that only the variable \texttt{factor} inside \texttt{mpow} may be lowered,
because it holds small exact integers.

Parsimony reaches the same conclusion independently: at a relative tolerance of
$10^{-8}$, one of eleven variables is lowered, and it is \texttt{factor}. Because
the ladder extends below \texttt{float}, the result is sharper than the original
tool could express --- \texttt{factor} holds $6$, $4$ and $1$, so it descends to an
8-bit format with \emph{zero} error.

We note a discrepancy for the record. The paper's stated all-single result is
$0.4999990927$; hardware \texttt{float32} evaluation of the program as published
gives $0.4999980926513672$, which parsimony's emulator matches bit for bit. The
trailing digits agree and only the sixth decimal differs, which is consistent
with a transcription error.
\todo{Confirm against the authoritative PDF text rather than a rendered page, and
if it stands, note it as an erratum rather than a disagreement.}

\subsection{NAS \texttt{cg}}
\label{sec:nascg}

\pending{The credibility anchor, and currently absent. NAS \texttt{cg} appears in
Precimonious' own benchmark set, so its published configurations (Table~4 of
\cite{blame}: 32 double variables initially, reducing to splits of 2/30, 13/19,
16/16 and 16/16 at thresholds $10^{-4}$ through $10^{-10}$) and speedups
($7.7$--$7.9\%$) are directly comparable. It is also one of only five programs in
that study showing any speedup, so it exercises both halves of the dual
objective.

Three checks, in order: (i) does the search find a comparable double/float split
at each threshold --- exact agreement is not expected, since the variable set at
bitcode level differs from a source-level set and the original was tuned at
\texttt{-O0}; (ii) what is the evaluation count, which is the baseline any
prefilter must beat; (iii) is threshold monotonicity respected, since a tighter
threshold admitting more lowering indicates a bug.}

\section{Case Study: Sparse Iterative Solvers}
\label{sec:petsc}

\pending{The result that would make this a paper rather than an architecture
note. Two questions in PETSc~\cite{petsc}, chosen because they are common targets
rather than exotic ones, and because both have a real payoff on current hardware:

\emph{How many bits does a Jacobian need?} Inexact Newton methods tolerate
inaccurate Jacobians by construction --- the outer iteration corrects them --- and
practitioners know this qualitatively. Nobody has published a number. This is a
direct application of virtual-width search to a quantity whose precision
requirement is genuinely unknown.

\emph{Mixed-precision preconditioning and matrix storage.} Sparse matrix-vector
products are bandwidth-bound, so storing the operator in \texttt{fp32} or
\texttt{bfloat16} while accumulating in \texttt{fp64} is close to a factor of two.
This is precisely the regime in which the original evaluation's $0.0\%$ becomes a
real number, which makes it the sharpest available test of this paper's thesis.

Both require an evaluator that can run compiled MPI code, which does not yet
exist (\S\ref{sec:limitations}).}

\section{Limitations and Future Work}
\label{sec:limitations}

\paragraph{Search cost.} Delta debugging costs $O(n^2)$ evaluations per ladder
rung. In the reproduction of \S\ref{sec:repro}, eleven variables at a tight
tolerance required 350 evaluations. Blame Analysis' prefilter is the established
remedy and is the most valuable unbuilt component.

\paragraph{No performance measurement.} The numpy evaluator answers how many bits
the arithmetic requires, never whether a configuration is faster. The dual
objective that motivates this paper requires an evaluator that compiles and
times real code. Until one exists, every performance claim here is an argument
about hardware trends rather than a measurement.

\paragraph{Granularity at scale.} Instruction-level attribution cost $50\times$
runtime and, in its offline form, 256\,GB on a small benchmark. Applying it to a
production library is not plausible. We believe the tractable variant is
attribution at the granularity of \emph{library operations} --- a sparse
matrix-vector product, an inner product --- rather than LLVM instructions. This
trades resolution for feasibility, and it shifts the hard step to deciding which
operations to attribute, which is a modeling judgment informed by reading the
source.

\paragraph{The optimization gap.} Like its predecessors, any source- or
bitcode-level variable search tunes a program that the optimizer will
subsequently transform. This remains unsolved and is a threat to the validity of
any configuration reported by tools in this family.

\paragraph{Accelerator code.} LLVM-instrumentation evaluators do not observe GPU
kernels, which is where the reduced-precision payoff is largest. Conclusions must
presently be established on a CPU path and confirmed empirically on the
accelerator.

\section{Related Work}

\paragraph{Search-based tuning.} Precimonious~\cite{precimonious} and Blame
Analysis~\cite{blame} are the direct ancestors. HiFPTuner~\cite{hifptuner} adds
static analysis to group intermediates and search hierarchically.
TAFFO~\cite{taffo} targets fixed-point conversion. POP~\cite{pop} takes a static,
constraint-based route.

\paragraph{Parallel revival.} PROMISE has recently been extended with FloatX to
support \texttt{bfloat16} and both 8-bit encodings, using CADNA's discrete
stochastic arithmetic for validation~\cite{promise2026}. That work reaches the
same conclusion as ours about the value of revisiting these methods for modern
formats, from an independent lineage, and reports that $43\%$ of variables admit
32 bits or fewer and $28\%$ admit 16 or fewer. Its benchmarks remain small and
serial; the gap at HPC scale is what motivates the present design.

\paragraph{Error estimation.} Verificarlo~\cite{verificarlo} instruments LLVM IR
to provide Monte Carlo arithmetic and variable-precision backends;
Verrou~\cite{verrou} achieves comparable ends via Valgrind without recompilation;
CADNA~\cite{cadna} implements discrete stochastic arithmetic through instrumented
types. All three are candidate evaluators under the protocol of
\S\ref{sec:design}.

\paragraph{Mixed precision by analysis rather than search.} Where a numerical
method's error is understood analytically, precision can be assigned by theory
rather than by search --- iterative refinement in three precisions being the
canonical example~\cite{carsonhigham}. The two approaches are complementary:
theory covers the kernel, search covers the surrounding code that no one has
analyzed.

\section{Conclusion}

\pending{Write once \S\ref{sec:nascg} and \S\ref{sec:petsc} carry results. The
conclusion should state what the measured payoff is on a bandwidth-bound solver
and whether it confirms the thesis of \S\ref{sec:intro}; if it does not, that is
also a publishable result and the paper should say so.}

\section*{Availability}

Source, tests and the case-study harnesses are at
\url{https://github.com/markadams/parsimony}, BSD-3-Clause. No code from
Precimonious (BSD-3-Clause, \copyright~2013 UC Berkeley) is reproduced; its
algorithms are reimplemented from their published descriptions.

\section*{Acknowledgments}

\todo{DOE/ASCR funding statement.}

\bibliographystyle{plainnat}
\bibliography{refs}

\end{document}
